\chapter{Rappresentazione caratteri}
\section{Idea}
Possiamo associare a ogni carattere un numero.

I caratteri possono essere rappresentati in:
\begin{itemize}
    \item \textbf{ASCII standard}: 1 carattere è rappresentato con 7 bit per un totale di 128 simboli rappresentabili
    \item \textbf{ASCII estesa}: 1 carattere è rappresentato con 8 bit rappresentabili fino a 256 simboli
    \item \textbf{UNICODE}: 1 carattere è rappresentato con un numero maggiore di bit (tra 8 e 32 bit per carattere)
\end{itemize}

\section{ASCII Standard}
\subsection{Caratteristiche}
\begin{itemize}
    \item ASCII standard contiene:
    \begin{itemize}
        \item 26 + 26 lettere (maiuscole + minuscole)
        \item 10 cifre decimali
        \item segni di interpunzione
        \item caratteri di controllo
    \end{itemize}

    \item Le cifre sono ordinate per valore
    \item Le lettere maiuscolo sono ordinate alfabeticamente
    \item Le lettere minuscole sono ordinate alfabeticamente
\end{itemize}

\newpage
\subsection{Tabella}
\begin{table}[htbp]
\centering
\small
\begin{tabular}{rrc|rrc|rrc|rrc}
\toprule
\textbf{Dec} & \textbf{Hex} & \textbf{Char} & \textbf{Dec} & \textbf{Hex} & \textbf{Char} & \textbf{Dec} & \textbf{Hex} & \textbf{Char} & \textbf{Dec} & \textbf{Hex} & \textbf{Char} \\
\midrule
0  & 00 & NUL & 32 & 20 & SPC & 64  & 40 & @  & 96  & 60 & ` \\
1  & 01 & SOH & 33 & 21 & !   & 65  & 41 & A  & 97  & 61 & a \\
2  & 02 & STX & 34 & 22 & "   & 66  & 42 & B  & 98  & 62 & b \\
3  & 03 & ETX & 35 & 23 & \#  & 67  & 43 & C  & 99  & 63 & c \\
4  & 04 & EOT & 36 & 24 & \$  & 68  & 44 & D  & 100 & 64 & d \\
5  & 05 & ENQ & 37 & 25 & \%  & 69  & 45 & E  & 101 & 65 & e \\
6  & 06 & ACK & 38 & 26 & \&  & 70  & 46 & F  & 102 & 66 & f \\
7  & 07 & BEL & 39 & 27 & '   & 71  & 47 & G  & 103 & 67 & g \\
8  & 08 & BS  & 40 & 28 & (   & 72  & 48 & H  & 104 & 68 & h \\
9  & 09 & HT  & 41 & 29 & )   & 73  & 49 & I  & 105 & 69 & i \\
10 & 0A & LF  & 42 & 2A & *   & 74  & 4A & J  & 106 & 6A & j \\
11 & 0B & VT  & 43 & 2B & +   & 75  & 4B & K  & 107 & 6B & k \\
12 & 0C & FF  & 44 & 2C & ,   & 76  & 4C & L  & 108 & 6C & l \\
13 & 0D & CR  & 45 & 2D & -   & 77  & 4D & M  & 109 & 6D & m \\
14 & 0E & SO  & 46 & 2E & .   & 78  & 4E & N  & 110 & 6E & n \\
15 & 0F & SI  & 47 & 2F & /   & 79  & 4F & O  & 111 & 6F & o \\
16 & 10 & DLE & 48 & 30 & 0   & 80  & 50 & P  & 112 & 70 & p \\
17 & 11 & DC1 & 49 & 31 & 1   & 81  & 51 & Q  & 113 & 71 & q \\
18 & 12 & DC2 & 50 & 32 & 2   & 82  & 52 & R  & 114 & 72 & r \\
19 & 13 & DC3 & 51 & 33 & 3   & 83  & 53 & S  & 115 & 73 & s \\
20 & 14 & DC4 & 52 & 34 & 4   & 84  & 54 & T  & 116 & 74 & t \\
21 & 15 & NAK & 53 & 35 & 5   & 85  & 55 & U  & 117 & 75 & u \\
22 & 16 & SYN & 54 & 36 & 6   & 86  & 56 & V  & 118 & 76 & v \\
23 & 17 & ETB & 55 & 37 & 7   & 87  & 57 & W  & 119 & 77 & w \\
24 & 18 & CAN & 56 & 38 & 8   & 88  & 58 & X  & 120 & 78 & x \\
25 & 19 & EM  & 57 & 39 & 9   & 89  & 59 & Y  & 121 & 79 & y \\
26 & 1A & SUB & 58 & 3A & :   & 90  & 5A & Z  & 122 & 7A & z \\
27 & 1B & ESC & 59 & 3B & ;   & 91  & 5B & [  & 123 & 7B & \{ \\
28 & 1C & FS  & 60 & 3C & \textless{}  & 92 & 5C & \textbackslash{} & 124 & 7C & \textbar{} \\
29 & 1D & GS  & 61 & 3D & =   & 93  & 5D & ]  & 125 & 7D & \} \\
30 & 1E & RS  & 62 & 3E & \textgreater{} & 94 & 5E & \textasciicircum{} & 126 & 7E & \textasciitilde{} \\
31 & 1F & US  & 63 & 3F & ?   & 95  & 5F & \_ & 127 & 7F & DEL \\
\bottomrule
\end{tabular}
\end{table}

\newpage
\section{ASCII Esteso}
\subsection{Caratteristiche}
Uguale a l'ASCII standard ma con più caratteri

\subsection{Tabella }
\begin{table}[htbp]
\centering
\small
\begin{tabular}{rrc|rrc|rrc|rrc}
\toprule
\textbf{Dec} & \textbf{Hex} & \textbf{Char} & \textbf{Dec} & \textbf{Hex} & \textbf{Char} & \textbf{Dec} & \textbf{Hex} & \textbf{Char} & \textbf{Dec} & \textbf{Hex} & \textbf{Char} \\
\midrule
128 & 80 & -- & 160 & A0 & NBSP & 192 & C0 & À & 224 & E0 & à \\
129 & 81 & -- & 161 & A1 & ¡ & 193 & C1 & Á & 225 & E1 & á \\
130 & 82 & -- & 162 & A2 & ¢ & 194 & C2 & Â & 226 & E2 & â \\
131 & 83 & -- & 163 & A3 & £ & 195 & C3 & Ã & 227 & E3 & ã \\
132 & 84 & -- & 164 & A4 & ¤ & 196 & C4 & Ä & 228 & E4 & ä \\
133 & 85 & -- & 165 & A5 & ¥ & 197 & C5 & Å & 229 & E5 & å \\
134 & 86 & -- & 166 & A6 & ¦ & 198 & C6 & Æ & 230 & E6 & æ \\
135 & 87 & -- & 167 & A7 & § & 199 & C7 & Ç & 231 & E7 & ç \\
136 & 88 & -- & 168 & A8 & ¨ & 200 & C8 & È & 232 & E8 & è \\
137 & 89 & -- & 169 & A9 & © & 201 & C9 & É & 233 & E9 & é \\
138 & 8A & -- & 170 & AA & ª & 202 & CA & Ê & 234 & EA & ê \\
139 & 8B & -- & 171 & AB & « & 203 & CB & Ë & 235 & EB & ë \\
140 & 8C & -- & 172 & AC & ¬ & 204 & CC & Ì & 236 & EC & ì \\
141 & 8D & -- & 173 & AD & SHY & 205 & CD & Í & 237 & ED & í \\
142 & 8E & -- & 174 & AE & ® & 206 & CE & Î & 238 & EE & î \\
143 & 8F & -- & 175 & AF & ¯ & 207 & CF & Ï & 239 & EF & ï \\
144 & 90 & -- & 176 & B0 & ° & 208 & D0 & Ð & 240 & F0 & ð \\
145 & 91 & -- & 177 & B1 & ± & 209 & D1 & Ñ & 241 & F1 & ñ \\
146 & 92 & -- & 178 & B2 & ² & 210 & D2 & Ò & 242 & F2 & ò \\
147 & 93 & -- & 179 & B3 & ³ & 211 & D3 & Ó & 243 & F3 & ó \\
148 & 94 & -- & 180 & B4 & ´ & 212 & D4 & Ô & 244 & F4 & ô \\
149 & 95 & -- & 181 & B5 & µ & 213 & D5 & Õ & 245 & F5 & õ \\
150 & 96 & -- & 182 & B6 & ¶ & 214 & D6 & Ö & 246 & F6 & ö \\
151 & 97 & -- & 183 & B7 & · & 215 & D7 & × & 247 & F7 & ÷ \\
152 & 98 & -- & 184 & B8 & ¸ & 216 & D8 & Ø & 248 & F8 & ø \\
153 & 99 & -- & 185 & B9 & ¹ & 217 & D9 & Ù & 249 & F9 & ù \\
154 & 9A & -- & 186 & BA & º & 218 & DA & Ú & 250 & FA & ú \\
155 & 9B & -- & 187 & BB & » & 219 & DB & Û & 251 & FB & û \\
156 & 9C & -- & 188 & BC & ¼ & 220 & DC & Ü & 252 & FC & ü \\
157 & 9D & -- & 189 & BD & ½ & 221 & DD & Ý & 253 & FD & ý \\
158 & 9E & -- & 190 & BE & ¾ & 222 & DE & Þ & 254 & FE & þ \\
159 & 9F & -- & 191 & BF & ¿ & 223 & DF & ß & 255 & FF & ÿ \\
\bottomrule
\end{tabular}
\end{table}

\newpage
\section{UNICODE}
\subsection{Caratteristiche}
\begin{itemize}
    \item È una evoluzione dello standard ASCII
    \item UNICODE è uno standard per la rappresentazione di caratteri
    \item Codifica tutti i caratteri utilizzati nelle principali lingue del mondo
    \item Indipendente dalla lingua, dal sistema operativo e dal programma utilizzato
    \item Inizalmente rappresentato come una codifica su 16 bit, ma poi esteso a 24 e 32 bit
    \item Unicode è in continua evoluzione e continua ad aggiungere sempre più caratteri
    \item Un carattere UNICODE è caratterizzato dal suo codice numerico, detto code point, solitamente rappresentato con 8 cifre esadecimali
    \item Con UNICODE, è possibile creare e gestire senza troppa pena documenti multilingue
    \item In particolare, tuttigli gli standard W3C supportano UNICODE
\end{itemize}
